\documentclass[11pt]{article}
\usepackage[margin=1in]{geometry}
\usepackage{amsmath,amssymb,amsthm,mathtools}
\usepackage{array,booktabs,longtable}
\usepackage[hyphens]{url}
\usepackage[hidelinks]{hyperref}
\newcommand{\Q}{\mathbb Q}
\newcommand{\F}{\mathbb F}
\newcommand{\SO}{\operatorname{SO}}
\newcommand{\tr}{\operatorname{tr}}
\newcommand{\adj}{\operatorname{adj}}
\newtheorem{theorem}{Theorem}[section]
\newtheorem{proposition}[theorem]{Proposition}
\newtheorem{corollary}[theorem]{Corollary}
\newtheorem{lemma}[theorem]{Lemma}
\theoremstyle{remark}
\newtheorem{remark}[theorem]{Remark}
\setlength{\emergencystretch}{3em}
\title{Issues and Resolutions to Open Problems in\\
  \emph{Rational Rotations of Three Dimensional Minkowski Space}}
\author{}
\date{12 October 2026}

\begin{document}
\maketitle
\begin{abstract}
This companion addresses N.~J.~Wildberger's revised Minkowski paper.
It records the constructive
formalisation of its finite algebraic assertions, clarifies null-axis spread
and the inverse of a negative-norm ray, and separates causal character from
arithmetic splitting. It answers the coefficient-product
question in Problem~90 and an explicit projective boost--rotation
factorization for Problem~92, including the half spread and spinor norm of
the Thomas--Wigner rotation. The general ternary Clifford correspondence
answers Problem~93. Finite-word spinor calculations give partial information
for Problem~91, while the integral programme in Problem~94 remains open.
Results with existing Lean certificates are distinguished from the
additional coordinate proofs supplied in this companion.
\end{abstract}

\section{Source, conventions and status}

The source is the supplied 36-page revision of \emph{Rational rotations of
three dimensional Minkowski space: rotations, boosts and spinors}
\cite{source}. Its numbered assertion inventory ends at Theorem~89 and its
research questions are Problems~90--94. Vectors are rows in the time-first
convention
\[
 G=\operatorname{diag}(1,-1,-1),\quad
 B(v,u)=vGu^T,\quad Q(t,r,s)=t^2-r^2-s^2.
\]
All general-field statements assume characteristic different from two.
The signed cross product is exactly the source's
\[
 v\times u=(s_vr_u-r_vs_u,\ s_vt_u-t_vs_u,\ t_vr_u-r_vt_u).
\]
The dihedron multiplication, norm and matrix map are
\begin{equation}
 (x:v)(y:u)=(xy-B(v,u):xu+yv-v\times u),\qquad
 N(x:v)=x^2+Q(v),
 \label{eq:product}
\end{equation}
\begin{equation}
 \Phi(x:t,r,s)=
 \begin{pmatrix}x+s&r-t\\r+t&x-s\end{pmatrix},
 \qquad\det\Phi=N.
 \label{eq:matrix}
\end{equation}
The product is multiplicative in this matrix order. The rotation invariant
is \(p(\rho)=(3-\tr\rho)/4\). For a cap represented by \((1:v)\),
\[
 p=\frac{Q(v)}{1+Q(v)},\qquad \chi=[1+Q(v)]=[1-p].
\]

The repository's older rational Minkowski development uses the different
form \(\operatorname{diag}(1,1,-1)\). Its vectors and cross-product signs
cannot be read as the coordinates above. The formal comparison uses an
explicit time-first form owner and a vector sign conversion; all examples
below use \eqref{eq:matrix} directly.

\begin{center}
\begin{tabular}{>{\raggedright\arraybackslash}p{0.17\textwidth}>{\raggedright\arraybackslash}p{0.74\textwidth}}
\toprule
Source item & Result of this companion \\
\midrule
Definition~45 / Theorem~46 & Extension to zero and nonzero null anchors,
with an explicit non-null test direction. \\
Section~12.3 & Exact inverse ray for both signs of the norm; distinction
between representative rays and anchor-vector rays. \\
Problem~90 & Necessity and sufficiency in characteristic different from two of
\(\delta=-\gamma^2\), by the coordinate proof below. \\
Problem~91 & Finite-word determinant classes and kernel tests; no complete
classification of generated subgroups or arithmetic subgroups. \\
Problem~92 & Projective Wigner factor, its half spread and spinor norm,
for two future pure boosts; factor order is specified. \\
Problem~93 & General ternary projective Clifford correspondence. \\
Problem~94 & Elementary primitive-representative reduction is described;
the integral classification programme is not solved. \\
\bottomrule
\end{tabular}
\end{center}

\section{Null axes and the measurement of spread}

Definition~45 defines spread for a cap with a non-null anchor, or a half turn.
The prose after Theorem~46 also discusses parabolic rotations with \(p=0\).
To attach its divided vector-spread measurement to these rotations, the
definition and existence argument should explicitly include null anchors.
The nondegenerate-perpendicular-plane argument in its printed proof alone
does not establish that extension.

For a nonzero null vector \(v=(t,r,s)\), a non-null perpendicular vector
can be chosen without a square root. If \(t\ne0\), take
\[
 w=(0,s,-r),\qquad B(v,w)=0,\quad Q(w)=-r^2-s^2=-t^2\ne0.
\]
If \(t=0\), take \(w=(1,0,0)\). The same latter vector handles the zero
anchor. For a non-null anchor the existing finite perpendicular construction
applies. Thus an admissible \(w\) exists in every case, over any field of
characteristic different from two.

The proof of the formula does not require \(Q(v)\ne0\). For any defined cap
and any supplied \(w\perp v\) with \(Q(w)\ne0\),
\[
 w\rho=\frac{(1-Q(v))w+2(v\times w)}{1+Q(v)},\qquad
 s(w,w\rho)=4p(1-p).
\]
Consequently a null cap has measured spread zero even when it is not the
identity. One proposed editorial repair is to let Definition~45 use any cap
anchor and insert the displayed null-anchor construction in Theorem~46.
This strengthens the stated domain without altering its formula. Non-null
test vectors remain essential: the denominator of vector spread is not
defined on a null test vector.

\section{Ray inverses and rational normalization}

\subsection{The sign-sensitive inverse}

A ray identifies \(P\) with \(aP\) only for \(a>0\). A projective
proportion allows every \(a\ne0\), so a non-null projective point contains
two opposite rays. The exact dihedron inverse is
\begin{equation}
 P^{-1}=\frac{\bar P}{N(P)},\qquad \overline{(x:v)}=(x:-v).
 \label{eq:inverse}
\end{equation}
This identity proves both inverse laws, since
\(P\bar P=\bar PP=N(P)1\). It also shows that inversion respects positive
scaling: \((aP)^{-1}=a^{-1}P^{-1}\).

For \(N(P)>0\), the inverse ray is the ray of \(\bar P\). For
\(N(P)<0\), it is the ray of \(-\bar P\). Using the conjugate alone in
that second case gives a product in the ray of \(-1\), not the identity
ray. Projectivization hides this sign, but positive rays retain it. The
source's full ray-group assertion is correct with \eqref{eq:inverse};
making the representative of the inverse explicit avoids this ambiguity.

For the source's negative-norm example,
\[
 P=(2:1,2,3),\quad N(P)=-8,\quad
 \Phi(P)=\begin{pmatrix}5&1\\3&-1\end{pmatrix},
\]
the inverse in the source coordinates is
\[
 P^{-1}=(-\tfrac14:\tfrac18,\tfrac14,\tfrac38),\qquad
 \Phi(P^{-1})=\begin{pmatrix}1/8&1/8\\3/8&-5/8\end{pmatrix}.
\]
Both matrix products are the identity; its norm is \(-1/8\). This exact
example and its sign conversion into the retained Clifford coordinates are
checked by the facade consumer.

\subsection{Two different sorts of rays}

Opposite dihedron rays \(P\) and \(-P\) represent the same rotation.
Faithfulness of the projective action says that its rotation fibers consist
of exactly these two rays, and the kernel is the scalar rays of \(1\) and
\(-1\). By contrast, opposite nonzero \emph{anchor-vector} rays \(v\)
and \(-v\) represent inverse cap rotations. This distinction holds for
timelike, spacelike and null anchors. It is the anchor-vector distinction
which records the source's sense of orientation.

The norm-one representative question is separate again. A rational ray can
be rescaled to norm one precisely when \(N(P)\) is a nonzero rational square.
If it exists, the positive rescaling and the unit representative are unique
within that ray. Negative norm makes normalization to \(+1\) impossible;
positive norm alone is insufficient over \(\Q\). The introductory boost
\((3:0,0,1)\) has norm eight and class \([2]\), so has no rational
determinant-one representative. Its projective matrix nevertheless gives
the exact rational boost
\[
 \begin{pmatrix}5/4&3/4&0\\3/4&5/4&0\\0&0&1\end{pmatrix}.
\]
Hence rational dihedron rays cover the whole group twice, whereas rational
determinant-one representatives cover only \(\ker\chi\).

\section{Causal type and arithmetic splitting}

Theorem~89 uses the discriminant
\[
 \Delta=-16p(1-p).
\]
For a prescribed non-null axis \(v\), write its cap anchor as \(\lambda v\),
where \(\lambda\ne0\). The existence equation is
\(\lambda^2Q(v)=p/(1-p)\), and therefore
\[
 \Delta=\bigl(4\lambda(1-p)\bigr)^2(-Q(v)).
\]
Thus the arithmetic splitting condition is precisely that \(-Q(v)\) be
a square in the ground field. It is not just that the axis is spacelike.
The retained square witness constructs both nontrivial eigenvalues and
actual nonzero eigenvectors of the original three-by-three matrix.

The distinction is already correct in the revised source; it is worth making
its exact examples explicit. The introductory boost has \(p=-1/8\) and
nontrivial eigenvalues \(2,1/2\). The rational axis \((1,2,3)\) has
\(Q=-12\), and is spacelike, but 12 has no rational square root. The
formal example with anchor \((1/4)(1,2,3)\) has \(p=-3\), time entry
\(15/2>0\), and no nontrivial rational eigenvalue. It is causally
boost-type while arithmetically nonsplit over \(\Q\).

The non-orthochronous example \((2:1,2,3)\) has \(p=3/2\), time entry
\(-9/4\), and characteristic polynomial
\[
 (X-1)\bigl((X+2)^2-3\bigr).
\]
Over a field supplied with \(r^2=3\), its other roots are \(-2\pm r\).
They are not rational. Causal sign ranges, determinant square classes and
ground-field eigenvalue splitting are therefore three separate pieces of
information; none substitutes for the other.

\section{Problem 90: associative products}

The revised question asks about
\[
 v\star u=\frac{v+u+\gamma(v\times u)}{1+\delta B(v,u)},
 \qquad\gamma\ne0.
\]

\begin{theorem}[Lorentz coefficient classification]
This product is associative wherever all four products in the two
bracketings are defined if and only if \(\delta=-\gamma^2\).
\end{theorem}
\begin{proof}
Take \(e_0=(1,0,0)\) and \(e_1=(0,1,0)\). For any supplied nonzero
\(a\) with \(D=1+\delta a^2\ne0\), the Lorentz coordinate
identities
\[
 B(e_0,e_0)=1,\quad B(e_0,e_1)=0,\quad
 e_0\times e_1=(0,0,1),\quad e_0\times(0,0,1)=-e_1
\]
give, for the triple \((ae_0,ae_0,e_1)\), the same two middle coordinates
as in the Euclidean probe:
\[
 \bigl((ae_0\star ae_0)\star e_1\bigr)_r=1,\qquad
 \bigl(ae_0\star(ae_0\star e_1)\bigr)_r
     =\frac{1-\gamma^2a^2}{1+\delta a^2}.
\]
All four denominators are nonzero. Equality forces
\((\delta+\gamma^2)a^2=0\). Try \(a=1\). If its denominator is zero,
then \(\delta=-1\); when \(\gamma^2=1\) the required relation is already
true, and otherwise the nonzero probe \(a=\gamma\) has denominator
\(1-\gamma^2\ne0\). These explicit alternatives prove necessity in every
stated characteristic.

For sufficiency the dilation \(v\mapsto\gamma v\) identifies the product
with \((v+u+v\times u)/(1-B(v,u))\), the opposite dihedron affine product.
The associative homogeneous matrix product \eqref{eq:product} proves
equality on every common affine domain. The dilation is invertible and
preserves precisely the displayed denominators.
\end{proof}

This answers the revised question over every field of characteristic
different from two, including its primary rational case. The rational
coefficient classification already has a Euclidean Lean certificate. The
Lorentz necessity calculation above is a proof supplied in this paper, not
a newly added standalone Lean theorem.

\section{Problem 91: finite spinor calculations and remaining subgroup theory}

The determinant class of a supplied rational invertible matrix is computable
exactly. A nonzero rational \(c\) has the square-class signature consisting
of its sign and the primes occurring to odd exponent in its reduced
numerator and denominator. Sign and prime parities add modulo two under
multiplication. The sign coordinate is essential in the Lorentz case:
\(\chi\) takes values in all rational nonzero square classes, not just the
positive ones.

Given finitely many generator matrices, every supplied finite signed word
can therefore be evaluated by exact two-by-two products. Multiplicativity
of the determinant gives its spinor class, and a word is in \(\ker\chi\)
precisely when its determinant is a rational square. Comparing the class of
one word times the inverse of another gives an exact kernel-coset test.
The finitely many generator sign/parity vectors span a finite binary space,
whose vectors and generator words can be enumerated. These calculations
describe the spinor image detected by the generators. They do not determine
the internal relations of the generated subgroup.

In particular, if the generators are already in \(\ker\chi\), every word
has trivial spinor norm, but it may still be a different group element.
There is no internal subgroup classification hidden in this test. The
repository's complete generated-subgroup certificate is for the Euclidean
positive-norm case; its exact rational matrix and square-class laws supply
the ingredients here, but no separate Lorentz subgroup-classification
certificate is asserted.

The complete Euclidean finite-order table in the other companion cannot be
silently transferred to the full Lorentz group. For example, a Lorentz half
turn has \(\chi=[Q(v)]\), which may be negative. The norm-square test answers
the question for any supplied matrix or order certificate, but a complete
classification of finite-order spinor values and natural arithmetic
subgroups is not supplied here. These parts of Problem~91 remain open.

\section{Problem 92: a projective Thomas--Wigner factor}

We specify the convention first: factor the product of two future pure
boosts as a future pure boost followed by a spatial rotation, in the source's
matrix order. Changing this convention changes the factor one calls the
Wigner rotation. No angle or square-root normalization is needed for the
following factorization.

\subsection{The factorization and its invariants}

Let the boost anchors be
\[
 v=(0,r_1,s_1),\qquad u=(0,r_2,s_2),\qquad
 r_i^2+s_i^2<1.
\]
These inequalities select the future pure-boost representatives
\[
 A=\Phi(1:v),\qquad B=\Phi(1:u),
\]
which are symmetric with positive determinant and positive trace. Define
\begin{equation}
 d=1-B(v,u)=1+r_1r_2+s_1s_2,\qquad
 k=r_1s_2-s_1r_2,\qquad H=d^2+k^2.
 \label{eq:wignerdata}
\end{equation}
Let \(r=r_1+r_2\), \(s=s_1+s_2\). The product dihedron is
\(P=(d:k,r,s)\), so \(AB=\Phi(P)\).

\begin{theorem}[Projective boost--rotation decomposition]
Over \(\Q\), or any ordered field, the quantities in
\eqref{eq:wignerdata} satisfy \(d>0\), \(H>0\). Put
\[
 R=\begin{pmatrix}d&-k\\k&d\end{pmatrix},\qquad
 S=\begin{pmatrix}
 H+sd-rk&rd+sk\\rd+sk&H-sd+rk
 \end{pmatrix}.
\]
Then
\begin{equation}
 AB=H^{-1}SR.
 \label{eq:factorization}
\end{equation}
The projective class of \(S\) is a future pure boost and that of \(R\)
is a spatial rotation. The corresponding Wigner rotation \(\omega\) has
\begin{equation}
 p(\omega)=\frac{k^2}{d^2+k^2},\qquad
 \chi(\omega)=[d^2+k^2].
 \label{eq:wignerinvariants}
\end{equation}
This spatial-rotation factor is unique among factorizations in this order
with a future pure-boost factor.
\end{theorem}
\begin{proof}
Write \(a=r_1r_2+s_1s_2\) and \(b_i=r_i^2+s_i^2<1\).
The identity \(a^2+k^2=b_1b_2\) gives \(a^2<1\), hence \(a>-1\)
and \(d=1+a>0\). Thus \(H>0\). Direct multiplication gives
\[
 \Phi(P)\adj(R)=S,\qquad R\adj(R)=HI,
\]
and therefore \eqref{eq:factorization}. The matrix \(S\) is symmetric and
\[
 \det S=(\det A)(\det B)H=(1-b_1)(1-b_2)H>0,
 \qquad\tr S=2H>0.
\]
Its dihedron has coordinates
\( (H:0,rd+sk,sd-rk)\). After scaling by \(H^{-1}>0\), its scalar
coordinate is one, its time anchor coordinate is zero, and its norm is
positive. Thus its spatial anchor has squared length less than one, exactly
the future pure-boost condition. The matrix \(R\) has dihedron coordinates
\((d:k,0,0)\), so it acts as a spatial rotation fixing the time axis.
The norm and half-spread identities give \eqref{eq:wignerinvariants}.

For uniqueness, any non-null spatial-rotation representative has the form
\(R'=\left(\begin{smallmatrix}a&-b\\b&a\end{smallmatrix}\right)\).
Symmetry of \(\Phi(P)\adj(R')\), required of its boost factor, is
equivalent to \(ka=db\). Since \(d\ne0\), nonzeroness of \(R'\) forces
\(a\ne0\), and its projective class is exactly that of \(R\).
\end{proof}

This gives a closed formula for the \emph{rotation factor}, rather than
merely detecting a time component in the product anchor. It stays over the
ground field: normalization of either two-by-two factor to determinant one
may require a square root, but the projective factors do not.

For comparison, the half spread and spinor norm of the \emph{whole product}
are
\begin{equation}
 p(\rho(v)\rho(u))
   =1-\frac{d^2}{(1+Q(v))(1+Q(u))},\qquad
 \chi(\rho(v)\rho(u))=[(1+Q(v))(1+Q(u))].
 \label{eq:wholeproduct}
\end{equation}
These are generally different from \eqref{eq:wignerinvariants}. They follow
directly from the homogeneous scalar coordinate and multiplicative norm;
calling them the Wigner rotation's invariants would conflate the boost with
the rotation factor.

\subsection{The perpendicular-boost example}

For Example~69 take
\[
 v=(0,0,1/3),\qquad u=(0,1/3,0).
\]
Then \(d=1\), \(k=-1/9\), and \(H=82/81\). Reversing the boosts
changes \(k\) to \(-k\) and inverts the Wigner rotation. The projective
rotation representative is \((1:-1/9,0,0)\), and
\[
 p(\omega)=\frac1{82},\qquad\chi(\omega)=[82].
\]
The whole product instead has \(p=-17/64\) and trivial spinor norm, since
both boost norms are \(8/9\). A convenient rational factorization is
\[
 \Phi(P)=
 \begin{pmatrix}4/3&4/9\\2/9&2/3\end{pmatrix}
 =\frac{81}{82}
 \begin{pmatrix}112/81&8/27\\8/27&52/81\end{pmatrix}
 \begin{pmatrix}1&1/9\\-1/9&1\end{pmatrix}.
\]
The residual future pure boost has anchor \((0,12/41,15/41)\). Its
spinor class is also \([82]\), so the two nontrivial factor classes cancel
in the product. This is an exact rational example of why the product's
spinor norm alone does not determine that of the Wigner rotation.

The projective product and Example~69 are already Lean-checked. The
factorization theorem, uniqueness argument and its Wigner formulas are
additional coordinate proofs in this paper. They have not been added as
standalone Lean declarations, and the earlier full audit should not be
represented as certifying those new theorem statements.

\section{Problem 93: the general ternary correspondence}

For a symmetric nonsingular matrix \(M\), put
\[
 B_M(u,v)=uMv^T,\quad Q_M(u)=B_M(u,u),\quad d=\det M,
 \quad C_M(u,v)=(u\times_{\mathrm{coord}}v)\adj(M).
\]
The general even Clifford product is
\[
 (x:u)(y:v)=(xy-dB_M(u,v):xv+yu-C_M(u,v)),
 \qquad N_M(x:u)=x^2+dQ_M(u).
\]
Its conjugation, associative product and multiplicative norm give the
action
\[
 (0:w)\longmapsto N_M(P)^{-1}\bar P(0:w)P.
\]
The retained constructive two-reflection theorem produces a preimage of
every supplied member of \(\SO(M,F)\), and equality of actions is equivalent
to an explicit nonzero scalar relation between the original non-null
representatives. This is the projective correspondence requested in
Problem~93, with the determinant factors and ownership of \(M\) retained.
The proof and affine formula are stated in the Euclidean companion
\cite{companion}; the same subject-owned Lean provider serves both papers.

For the present \(G\), \(d=1\) and the generic form-relative cross is the
\emph{negative} of the paper's signed Lorentz cross. Negating the vector
coordinate of the generic Clifford representative identifies its product
with \eqref{eq:product}. This explicit conversion explains why one may
share the abstract theorem without changing the source's signs.

The null quadric has a separate exact description. A nonzero null dihedron
maps to a nonzero singular two-by-two matrix. Testing its four possible
nonzero pivot entries constructs a nonzero column \(\xi\) and row \(\eta\)
with \(\Phi(P)=\xi\eta\). Its only factorization ambiguity is
\[
 (\xi,\eta)\longmapsto(a\xi,a^{-1}\eta),\qquad a\ne0.
\]
Holding either projective spinor fixed gives a ruling; distinct rulings in
one family are disjoint, while a column ruling and a row ruling meet exactly
in their outer-product point. The finite-field lists have exactly
\((q+1)^2\) such projective null points and \(q(q^2-1)\) non-null rotation
points. These counts and all four pivot charts are part of the checked
Minkowski facade surface, not consequences of an assumed completed quadric.

\section{Problem 94: the integral programme}

Every supplied rational non-null proportion can be reduced to an integral
primitive four-vector: clear the four denominators, divide by their common
integer gcd, and choose the sign of the first nonzero coordinate. This gives
a unique representative for its rational projective point. Two primitive
integer four-vectors representing the same rational point differ only by
sign: a reduced rational scale's denominator must divide all four input
coordinates, and its numerator must divide all four output coordinates.
Both gcds are one.

Products of integral dihedrons are integral by \eqref{eq:product}, and their
norms multiply. Thus one finite composition is performed by multiplying the
four coordinates and then removing the resulting common gcd and sign.
These elementary observations specify the arithmetic operation, but do not
give a classification of its primitive norm fibers or integral subgroups.

The matrix lattice deserves an explicit qualification. An integral matrix
\(\left(\begin{smallmatrix}a&b\\c&e\end{smallmatrix}\right)\)
has integral dihedron coordinates exactly when
\[
 a\equiv e\pmod2,\qquad b\equiv c\pmod2,
\]
because its inverse coordinate map is
\[
 x=(a+e)/2,\quad s=(a-e)/2,\quad
 r=(b+c)/2,\quad t=(c-b)/2.
\]
Primitive four-coordinates do not necessarily give a primitive matrix:
\((1:0,0,3)\) is primitive and non-null, but its matrix
\(\operatorname{diag}(4,-2)\) has gcd two. Conversely, a primitive matrix
need not lie in this parity lattice before rescaling. An integral theory must
specify which primitive object it normalizes; rational projective equality
does not by itself settle that lattice question.

These are introductory arithmetic reductions proved directly here. No
standalone Lean implementation of the primitive-normalization procedure,
full arithmetic orbit classification, presentation of integral subgroups or
local-to-global norm theorem is claimed. Problem~94 remains open as that
larger programme.

\section{Formal evidence, finite boundaries and open work}

The facade is
\path{formalisations/MinkowskiRotations/MinkowskiRotations.lean}, and its
claim ledger is \path{formalisations/MinkowskiRotations/coverage.md}.
It covers the finite algebraic assertions numbered 1--89 and the unnumbered
constructions, with direct source-coordinate facade consumers. Entry points
include the following owners.

\begingroup\small
\begin{longtable}{>{\raggedright\arraybackslash}p{0.28\textwidth}>{\raggedright\arraybackslash}p{0.65\textwidth}}
\toprule
Checked result & Mathematical owner \\
\midrule
\endhead
Source vector and anchor signs &
\path{Purelib/Geometry/Quadratic/Ternary/Lorentz.lean}. \\
Dihedron/matrix dictionary &
\path{Purelib/Algebra/Clifford/Lorentz/Matrix.lean}; exact product,
determinant and two coordinate inverses. \\
Composition and velocities &
\path{Purelib/Geometry/Rotations/Ternary/Lorentz/Composition.lean};
the partial anchor law and exact Einstein velocity identity. \\
Ray inversion and fibers &
\path{Purelib/Algebra/Clifford/Ternary/Ray.lean};
\path{Purelib/Geometry/Rotations/Ternary/Clifford/RayCover.lean}. \\
Prescribed axes and splitting &
\path{Purelib/Geometry/Rotations/Ternary/Lorentz/Prescribed.lean};
\path{Purelib/Geometry/Rotations/Ternary/Lorentz/Splitting.lean}. \\
Null spinors and rulings &
\path{Purelib/LinearAlgebra/Coordinates2/RankOne.lean};
the original Lorentz matrix recovery and finite null-point enumeration. \\
General ternary correspondence &
\path{Purelib/Geometry/Rotations/Ternary/Clifford.lean};
\texttt{projectiveCliffordCorrespondence}. \\
\bottomrule
\end{longtable}
\endgroup

The mathematical tree at commit \texttt{a6219f044} passed all twelve
completion gates on 11 October 2026, including the warning-fatal full build
and lint, full configured-script axiom audit and independent nanoda pass.
Build phases used 64 Lean threads. The normal audit checked 136,660 project
declarations and 1,536 selected canonical declarations with no allowed
axioms; nanoda accepted all 42 compatible module sets. This is the selected
mathematical proof closure, not a strict audit of every internal meta-provider
helper. The Lorentz proof for Problem~90, the Wigner factorization for
Problem~92 and the integral reductions for Problem~94 are printed proofs
introduced here, outside that existing Lean declaration inventory.

The infinite-index assertion of Corollary~86 has retained finite content:
for every supplied \(n\), construct \(n\) pairwise distinct rational square
classes, preimages with those norms and pairwise kernel-separated rotations.
The source norm witness
\[
 \left(\frac{c+1}{2}:0,\frac{c-1}{2},0\right)
 \quad\text{has norm }c
\]
works for every requested nonzero rational \(c\). This supplies arbitrary
finite coset packs without forming a completed quotient carrier. Real square
roots, topological connected components, completed Lie groups and analytic
angle/rapidity dictionaries remain outside the formal scope. Their absence
is a constructive boundary, not a claimed defect in the source.

Problem~90 has the general-field answer above. Problem~92 has the displayed
projective decomposition for the stated pure-boost convention, and
Problem~93 has the general constructive correspondence. Problem~91 still
asks for internal subgroup and arithmetic classifications not obtained by
finite spinor tests. Problem~94 still asks for an integral theory beyond
primitive-coordinate reduction. Formalising the additional printed proofs
would be separate work; these companions do not imply that it has already
been done.

\begin{thebibliography}{9}
\bibitem{source} N.~J.~Wildberger,
\emph{Rational rotations of three dimensional Minkowski space: rotations,
boosts and spinors}, supplied revision
\texttt{RelativisticsRationalRotations3.pdf}, 36 pages.
Source SHA-256:
\url{8c397b21f726bf0a8c89b8c56bf84edf5d2586bcac552c0f10134d3c9fff46d7}.
\bibitem{companion} \emph{Issues and Resolutions to Open Problems in Rational
Rotations of Three Dimensional Euclidean Space}, companion in
\path{formalisations/RationalRotations/IssuesAndOpenProblems.pdf},
12 October 2026.
\end{thebibliography}
\end{document}
